% Solutions of Book 4, chapter 3 (cpp-polymorphism).
\section*{Chapter \ref{ch:b2:cpp-polymorphism} --- Inheritance and Polymorphism}

\begin{solution}{exo:b2:cpp-polymorphism:1}
\omoutput{code/bachelor-2/03-cpp-polymorphism/outputs/exo_which.txt}{The four lines.}
\code{kind} is virtual: through the pointer and the reference it is dispatched to \code{Map::kind}. \code{shelf} is not
virtual: through an \code{Item *} or \code{Item &} the compiler calls \code{Item::shelf}, chosen at compile time from
the static type; \code{Map::shelf} only hides it, and is called when the static type is \code{Map}, as in the last
line. By value, the parameter is a sliced \code{Item}: both calls run \code{Item}'s versions.
\end{solution}

\begin{solution}{exo:b2:cpp-polymorphism:2}
\code{std::unique_ptr<Sensor>} destroys its object with \code{delete} on a \code{Sensor *}. \code{Sensor}'s
\omterm{def:b2:cpp-classes-raii:constructor}{destructor} is not virtual, so for the \code{LoggingSensor} only \code{~Sensor} runs (the vector of 1000 readings is
never destroyed) and the block is returned with the size of a \code{Sensor}, 4 bytes, instead of the 32 allocated:
undefined behaviour, which the sanitizer reports as a type mismatch, from inside \code{unique_ptr}'s \omterm{def:b2:cpp-classes-raii:constructor}{destructor}. The
fix is the word \code{virtual}: \code{virtual ~Sensor() {}}. (A \code{std::shared_ptr<Sensor>} made by
\code{make_shared<LoggingSensor>} would happen to work, because it remembers the real type; relying on that is
fragile.)
\end{solution}

\begin{solution}{exo:b2:cpp-polymorphism:3}
\omcode{code/bachelor-2/03-cpp-polymorphism/exercises/triangle.hpp}{7}{23}{\code{Triangle}, in its own file}
Step 1 is already done by \code{Shape}; the file is step 2: one \omterm{def:b2:cpp-polymorphism:inheritance}{derived class}, overriding the pure \code{area} (by
Heron's formula) and \code{name}. The loop of \code{shapes.cpp} is step 3 and is untouched; only the place that creates
the shapes (step 4) gains a line, a \code{push_back} of \code{std::make_unique<Triangle>(3, 4, 5)}. The strict
inequalities refuse flat triangles, whose area would be 0. Tested in \code{tests/shapes_test.cpp}.
\end{solution}

\begin{solution}{exo:b2:cpp-polymorphism:4}
\code{Plain}: one \code{int}, 4 bytes. \code{OneVirtual}: \omterm{def:b2:cpp-polymorphism:vtable}{table pointer} (8) and \code{int} (4) make 12, rounded up to
16, a multiple of the pointer's alignment 8. \code{ManyVirtual}: the same 16, since four \omterm{def:b2:cpp-polymorphism:virtual}{virtual functions} still mean
one pointer (\cref{prop:b2:cpp-polymorphism:layout}). \code{Derived}: its \code{int} goes into the 4 bytes of padding
at the end of the base part, which this compiler's rules allow for a \omterm{def:b2:cpp-polymorphism:inheritance}{base class} with \omterm{def:b2:cpp-polymorphism:virtual}{virtual functions}, so it stays
16. A \omterm{def:b2:cpp-polymorphism:vtable}{table pointer} and three \code{char}: $8 + 3 = 11$, rounded to 16. Two \code{double} and five \omterm{def:b2:cpp-polymorphism:virtual}{virtual functions}:
$8 + 16 = 24$, already a multiple of 8.
\end{solution}

\begin{solution}{exo:b2:cpp-polymorphism:5}
\omoutput{code/bachelor-2/03-cpp-polymorphism/outputs/exo_ctor_virtual.txt}{The four lines.}
While \code{Item}'s \omterm{def:b2:cpp-classes-raii:constructor}{constructor} runs, the object is an \code{Item}: its \omterm{def:b2:cpp-polymorphism:vtable}{table pointer} is \code{Item}'s and the call
runs \code{Item::describe}. Then \code{Map}'s \omterm{def:b2:cpp-classes-raii:constructor}{constructor} sets the \omterm{def:b2:cpp-polymorphism:vtable}{table pointer} to \code{Map}'s, builds
\code{region}, and its call runs \code{Map::describe}; so does \code{show}. When \code{m} dies, \code{~Map} (implicit)
runs first, then \code{~Item}, during which the object is an \code{Item} again. If the base \omterm{def:b2:cpp-classes-raii:constructor}{constructor} called
\code{Map::describe}, it would read \code{region} before its \omterm{def:b2:cpp-classes-raii:constructor}{constructor} ran: a string not yet built, which is
undefined behaviour and in practice a crash or garbage.
\end{solution}

\begin{solution}{exo:b2:cpp-polymorphism:6}
\code{std::vector<Item>} holds \code{Item} objects: \code{push_back(r)} copies the \code{Item} part of the recording
(slicing), so the element is a plain \code{Item} and its \code{label} is \code{Item}'s. To keep objects of several
classes, the vector must hold pointers that own them: \code{std::vector<std::unique_ptr<Item>>}, filled with
\code{std::make_unique<Recording>(...)}; the loop then calls \code{it->label()} and dispatch works. A vector of
references does not exist: a reference is not an object, cannot be reseated, and cannot be an element.
\code{std::unique_ptr} rather than \code{std::shared_ptr}, since the catalogue is the only owner.
\end{solution}

\begin{solution}{exo:b2:cpp-polymorphism:7}
\omcode{code/bachelor-2/03-cpp-polymorphism/exercises/fines_tables.h}{5}{27}{Fines with one table per kind, in C}
\code{tests/fines_tables_test.c} compares every kind and every lateness from $-2$ to 400 with the switch. To add a
kind: write its fine function and its table, two short declarations, possibly in a new file; nothing else changes.
The C version has no compiler check that a table is complete, which C++'s \omterm{def:b2:cpp-polymorphism:abstract}{pure virtual functions} give.
\end{solution}

\begin{solution}{exo:b2:cpp-polymorphism:8}
For \code{set_width(w)}: afterwards the width is \code{w} and the height is unchanged; symmetrically for
\code{set_height(h)}. From these, a client proves that after \code{set_width(4); set_height(5);} the area is
$4 \times 5 = 20$, whatever the starting sides: this is \code{resize_plot}. Suppose a \code{Square} keeps the invariant
width $=$ height. After \code{set_height(5)} the invariant forces the width to 5, so the area is 25, not 20: the
client's proven property fails, and the principle is violated. If instead the \omterm{def:b2:cpp-polymorphism:virtual}{override} left the width at 4, the
object would no longer be square, breaking its own invariant. No \omterm{def:b2:cpp-polymorphism:virtual}{override} can do both, since the two requirements
contradict each other on the state after \code{set_height(5)}. With immutable classes, \code{Rectangle(w, h)}
promises only that \code{area()} returns $wh$ for the sides given at construction; \code{Square(s)} constructs
\code{Rectangle(s, s)}, and every property provable from that specification (the area is the product of the sides,
and never changes) holds for it. The problem was mutation, not squares.
\end{solution}

\begin{solution}{exo:b2:cpp-polymorphism:9}
\omcode{code/bachelor-2/03-cpp-polymorphism/exercises/capped.hpp}{6}{19}{A policy that wraps a policy}
\code{CappedFine} derives from \code{FinePolicy}, so that the summing code can use it like any policy
(\omterm{def:b2:cpp-polymorphism:inheritance}{inheritance} of the interface); and it holds the policy it caps, to which it delegates (\omterm{def:b2:cpp-polymorphism:composition}{composition} of the
behaviour). It can therefore cap any policy, present or future, without a capped version of each. On the town's loans
at the end of 2025, the total owed falls from 463.45 to 309.05 with a cap of 10.00, and a cap above every fine gives
463.45 again; \code{tests/library_test.cpp} checks both figures.
\end{solution}

\begin{solution}{exo:b2:cpp-polymorphism:10}
\omcodes{code/bachelor-2/03-cpp-polymorphism/exercises/expr.hpp}{52-73,90-106}{The shared base of sums and products,
and products}
\code{Expr} declares \code{eval}, \code{str}, \code{simplify} (pure) and \code{constant}, which returns the address of
the value for a number and null otherwise; \code{Num} and \code{Var} are leaves. \code{Binary} writes \code{eval},
\code{str} and \code{simplify} once, in terms of three pure functions: the operation, its symbol, and
\code{rebuild}, which applies the identities of each operator to the simplified operands. \code{simplify} returns a new
tree and leaves the original alone. \code{tests/exercises_test.cpp} checks the example ($x = 1.5$ gives 7.5 before and
after), the identities, and an unbound variable.
\end{solution}

\begin{solution}{exo:b2:cpp-polymorphism:11}
Tables: $8 \sum_i m_i$ bytes, once for the program. \omterm{def:b2:cpp-polymorphism:vtable}{Table pointers}: $8 \sum_i N_i$ bytes. For the library: tables
$8 \times 3 \times 3 = 72$ bytes, three policy objects with one pointer each, 24 bytes, and the 400 loans' pointers
to their policy, 3200 bytes, against 400 tags of 4 bytes, 1600 bytes, in the C design: a few kilobytes either way.
For the game: tables $8 \times 10 \times 4 = 320$ bytes, negligible; each particle is $24 + 8 = 32$ bytes, so \omterm{def:b2:cpp-polymorphism:vtable}{table
pointers} are a quarter of the memory (8\,MB of 32\,MB). Memory per object, not code size, is the price of \omterm{def:b2:cpp-polymorphism:vtable}{dispatch
tables} when objects are small and many.
\end{solution}

\begin{solution}{exo:b2:cpp-polymorphism:12}
\omcode{code/bachelor-2/03-cpp-polymorphism/exercises/collide.hpp}{9}{42}{Double dispatch}
\code{a.hits(b)} is dispatched on \code{a}'s class: \code{Disc::hits} calls \code{b.hits_disc(*this)}, which is
dispatched on \code{b}'s class and knows that its argument is a \code{Disc}. A disc and a box intersect when the point
of the box nearest the disc's centre (the centre clamped to the box) is within the radius. With $c$ classes, each
needs \code{hits} and one \code{hits_X} per class: $c + 1$ functions each, $c(c+1)$ in all. Adding a class means
writing its $c + 2$ functions and adding one \code{hits_New} to every existing class: double dispatch is closed to new
kinds, the reverse of what \cref{met:b2:cpp-polymorphism:open-closed} achieves for single dispatch.
\end{solution}

\begin{solution}{pb:b2:cpp-polymorphism:1}
\textbf{1.} The \code{enum}, a case in \code{fine_cents}, a case in \code{kind_name}, and every item that should
become a map.

\textbf{2.} $fk$ cases (one per function and kind). Adding a kind writes $f$ cases, in $f$ different places, plus the
enum.

\textbf{3.} Both switches have no \code{default}, so \code{-Wall} (through \code{-Wswitch}) warns for both after
\code{MAP} is added; a switch with a \code{default} would stay silent.

\textbf{4.} $2 \times 50 = 100$ cents, 1.00; $10 \times 50 = 500$, capped to 4.00; and 4.00 for 60 days.

\textbf{5.} The kind arrives as data, a string, and some code must turn each string into the right class: a table or
a chain of tests is unavoidable at the boundary. Classes reduce the number of such places to one.

\textbf{6.} In \code{library.hpp}, in \code{policy_for}.

\textbf{7.} A new file holding \code{MapFine}, overriding \code{fine_cents} (\code{std::min(50 * d, 400)}) and
\code{name}; \code{policy_for} gains a line for the genre of maps. \code{FinePolicy}, the other policies and the
summing code do not change.

\textbf{8.} It now describes the rules of a kind, not only its fine: a name such as \code{KindRules} fits better. This
is step 1 of \cref{met:b2:cpp-polymorphism:open-closed}: collect what every kind must provide.

\textbf{9.} Again $fk$ function bodies, as many as cases, but grouped by kind, one class per kind, instead of by
function: a new kind adds its $f$ bodies in one place.

\textbf{10.} Comic: 19 loans, 51.75, about 2.72 each. Standard: 58 loans, 114.10, about 1.97 each. Study: 32 loans,
297.60, 9.30 each.

\textbf{11.} $6240 = 20(d - 3)$ gives $d = 315$. Loan 27 was due on 19 February 2025 and is still out: 315 days to 31
December.

\textbf{12.} Only \code{policy_for}, which returns a capped study policy; \code{CappedFine} itself is a new file.

\textbf{13.} Precondition: $d \geq 1$, a whole number of days late (the summing code calls it only for late loans).
Postcondition: the result is a whole number of cents, $\geq 0$, and the function does not throw. The summing code
needs nothing more.

\textbf{14.} Standard (at most 5.00) and comic (at most 3.00); study does not (11.40 for 60 days). It is not in the
specification of \code{FinePolicy}, only a property of two of its classes: the letter promises too much.

\textbf{15.} Yes: $\min(10d, 500)$ and $\min(25d, 300)$ are non-negative and non-decreasing, and
$\max(0, 20(d-3))$ too. \code{CappedFine} keeps both properties of the policy it wraps, since $\min(f(d), c)$ with $c
\geq 0$ is non-negative and non-decreasing when $f$ is.

\textbf{16.} Non-negativity (and the precondition: $d = 0$ is not a late loan). Summed with the fines, refunds would
lower the total ``owed'' and could make it negative; and since the summing code skips $d = 0$, the refund would never
be paid anyway. A refund is a different operation and belongs in a different function.

\textbf{17.} For a client relying on the 5.00 ceiling: no, a study loan 60 days late breaks it. For a client relying
only on question 13: yes. Substitutability is always relative to a specification.

\textbf{18.} No. \code{ComicFine} would inherit \code{StandardFine}'s promise of 10 cents a day and break it. What
they share, ``a rate per day with a ceiling'', is better expressed by one class with two parameters, or by
\code{CappedFine} around a linear policy.

\textbf{19.} In the printed run, a few nanoseconds more per call: the mixed-kinds line is about twice the switch
line. (The digits change from run to run; the ratio is what to read.)

\textbf{20.} 400 calls of a few nanoseconds: a few microseconds a night, nothing.

\textbf{21.} $10^8$ calls at the mixed-kinds rate, several nanoseconds each: most of each second, before any useful work.

\textbf{22.} Keep the particles of each kind in their own contiguous array and loop over each array with direct
calls (no pointer to follow, no unpredictable target), or keep one array of tagged values and switch; a template
(\cref{ch:b2:cpp-templates}) gives the first without writing each loop by hand.

\textbf{23.} Both cases do the same two loads and the same indirect call, and their objects are allocated the same
way. What differs is the predictability of the target: with one kind the processor guesses the next function
correctly every time, with a random mix half the time, and each wrong guess costs several nanoseconds
(\cref{ch:b2:processor-and-pipelining}).

\textbf{24.} \emph{One new file per new kind, and the measured cost of a virtual call against a switch}: with
classes, adding a kind is one new file and one line in the function that creates the objects, against $f$ cases in
$f$ places with a switch; a virtual call costs a few nanoseconds, in the printed run about twice an inlined switch
when the kinds are mixed, which matters only in inner loops over millions of objects.
\end{solution}
